\documentclass[10pt,a4paper]{amsart}
\usepackage{amsmath,amssymb,amsthm,graphicx}
% Retained minimal preamble for the Stacks Project chapter more-morphisms.tex (More on Morphisms).
% Verbatim excerpt of the chapter preamble: Xy-pic setup (native preamble lines 14-18) and the native
% theorem declarations with the mathematical macros (native preamble lines 154-217). Copied mechanically;
% nothing added, removed or redefined. The wrapper class and the amsmath/amssymb/amsthm/graphicx base
% packages are supplied by the assembler, so the class/externaldocument part of the native preamble is omitted.
% For commutative diagrams we use Xy-pic
\usepackage[all]{xy}

% We use 2cell for 2-commutative diagrams.
\xyoption{2cell}
% Theorem environments.
%
\theoremstyle{plain}
\newtheorem{theorem}[subsection]{Theorem}
\newtheorem{proposition}[subsection]{Proposition}
\newtheorem{lemma}[subsection]{Lemma}

\theoremstyle{definition}
\newtheorem{definition}[subsection]{Definition}
\newtheorem{example}[subsection]{Example}
\newtheorem{exercise}[subsection]{Exercise}
\newtheorem{situation}[subsection]{Situation}

\theoremstyle{remark}
\newtheorem{remark}[subsection]{Remark}
\newtheorem{remarks}[subsection]{Remarks}

\numberwithin{equation}{subsection}

% Macros
%
\def\lim{\mathop{\mathrm{lim}}\nolimits}
\def\colim{\mathop{\mathrm{colim}}\nolimits}
\def\Spec{\mathop{\mathrm{Spec}}}
\def\Hom{\mathop{\mathrm{Hom}}\nolimits}
\def\Ext{\mathop{\mathrm{Ext}}\nolimits}
\def\SheafHom{\mathop{\mathcal{H}\!\mathit{om}}\nolimits}
\def\SheafExt{\mathop{\mathcal{E}\!\mathit{xt}}\nolimits}
\def\Sch{\mathit{Sch}}
\def\Mor{\mathop{\mathrm{Mor}}\nolimits}
\def\Ob{\mathop{\mathrm{Ob}}\nolimits}
\def\Sh{\mathop{\mathit{Sh}}\nolimits}
\def\NL{\mathop{N\!L}\nolimits}
\def\CH{\mathop{\mathrm{CH}}\nolimits}
\def\proetale{{pro\text{-}\acute{e}tale}}
\def\etale{{\acute{e}tale}}
\def\QCoh{\mathit{QCoh}}
\def\Ker{\mathop{\mathrm{Ker}}}
\def\Im{\mathop{\mathrm{Im}}}
\def\Coker{\mathop{\mathrm{Coker}}}
\def\Coim{\mathop{\mathrm{Coim}}}

% Boxtimes
%
\DeclareMathSymbol{\boxtimes}{\mathbin}{AMSa}{"02}

%
% Macros for moduli stacks/spaces
%
\def\QCohstack{\mathcal{QC}\!\mathit{oh}}
\def\Cohstack{\mathcal{C}\!\mathit{oh}}
\def\Spacesstack{\mathcal{S}\!\mathit{paces}}
\def\Quotfunctor{\mathrm{Quot}}
\def\Hilbfunctor{\mathrm{Hilb}}
\def\Curvesstack{\mathcal{C}\!\mathit{urves}}
\def\Polarizedstack{\mathcal{P}\!\mathit{olarized}}
\def\Complexesstack{\mathcal{C}\!\mathit{omplexes}}
% \Pic is the operator that assigns to X its picard group, usage \Pic(X)
% \Picardstack_{X/B} denotes the Picard stack of X over B
% \Picardfunctor_{X/B} denotes the Picard functor of X over B
\def\Pic{\mathop{\mathrm{Pic}}\nolimits}
\def\Picardstack{\mathcal{P}\!\mathit{ic}}
\def\Picardfunctor{\mathrm{Pic}}
\def\Deformationcategory{\mathcal{D}\!\mathit{ef}}
\newenvironment{reference}{}{}
\newenvironment{slogan}{}{}
\newenvironment{history}{}{}
\usepackage{cleveref}
\crefname{originallemma}{original lemma}{original lemmas}
\Crefname{originallemma}{Original lemma}{Original lemmas}
\makeatletter
\newlabel{definition-pseudo-coherent}{{067Z}{source}{}{sourcepointer-1}{}}
\newlabel{definition-pseudo-coherent@cref}{{[originallemma][0][]067Z}{[1][0][]source}{}{}{}}
\makeatother
\crefname{originallemma}{original lemma}{original lemmas}
\Crefname{originallemma}{Original lemma}{Original lemmas}
\makeatletter
\newlabel{lemma-base-change-relative-pseudo-coherent}{{09UR}{source}{}{sourcepointer-2}{}}
\newlabel{lemma-base-change-relative-pseudo-coherent@cref}{{[originallemma][0][]09UR}{[1][0][]source}{}{}{}}
\makeatother
\crefname{originallemma}{original lemma}{original lemmas}
\Crefname{originallemma}{Original lemma}{Original lemmas}
\makeatletter
\newlabel{more-algebra-lemma-pull-relative-pseudo-coherent}{{067B}{source}{}{sourcepointer-3}{}}
\newlabel{more-algebra-lemma-pull-relative-pseudo-coherent@cref}{{[originallemma][0][]067B}{[1][0][]source}{}{}{}}
\makeatother
\crefname{originallemma}{original lemma}{original lemmas}
\Crefname{originallemma}{Original lemma}{Original lemmas}
\makeatletter
\newlabel{perfect-lemma-closed-push-pseudo-coherent}{{09VA}{source}{}{sourcepointer-4}{}}
\newlabel{perfect-lemma-closed-push-pseudo-coherent@cref}{{[originallemma][0][]09VA}{[1][0][]source}{}{}{}}
\makeatother
\title{Relative pseudo-coherence transfers between base and intermediate when the intermediate structure sheaf is pseudo-coherent}
\author{The Stacks Project Authors}
\date{stacks-project:more-morphisms.tex:local-snapshot-20261011; GNU FDL 1.2 or later}
\begin{document}
\maketitle
\noindent\textit{Editorial note.} Statement span is lines 17234-17246 of sources/stacks-more-morphisms/pinned-original-source.tex; proof span is lines 17248-17272, the complete author proof (23 body lines).\par\medskip
\noindent\textit{Editorial note.} Navigation stubs: 4 cited-but-not-included labels are supplied as non-load-bearing pointers with their pinned tags where the snapshot tag index has them.\par\medskip
\noindent\textit{Editorial note.} extra\_packages is empty on purpose: the retained preamble already loads xy with options and a second declaration clashes in the assembler wrapper.\par\medskip
\noindent\textit{Editorial note.} Text is retained verbatim, including the author's own typos and placeholders.\par\medskip
\noindent\textit{Editorial note.} Composition law of the chapter's relative-pseudo-coherence family; the footnote's dictionary (O\_Y pc over S means Y -> S pc) is retained.\par\medskip
\noindent\textit{Allowed standard prerequisites (editorial).}\begin{itemize}
\item Explicit prerequisite boundary. Definition of relative pseudo-coherence and pseudo-coherent morphisms (chapter definition section with the footnote dictionary); More on Algebra, lemma pull-relative-pseudo-coherent (the diagram technique the proof copies); chapter lemma base-change-relative-pseudo-coherent; Derived Categories of Schemes, lemma closed-push-pseudo-coherent.
\end{itemize}\medskip
\ifdefined\hypertarget\hypertarget{sourcepointer-1}{}\fi
\noindent\textit{Source reference.} Original lemma 067Z: sources/stacks-more-morphisms/pinned-original-source.tex, source lines 17375--17381. Cited by the retained statement or proof but not part of the collected spans; supplied only as a navigation pointer. The argument does not depend on its content.\par\medskip
\ifdefined\hypertarget\hypertarget{sourcepointer-2}{}\fi
\noindent\textit{Source reference.} Original lemma 09UR: sources/stacks-more-morphisms/pinned-original-source.tex, source lines 17112--17120. Cited by the retained statement or proof but not part of the collected spans; supplied only as a navigation pointer. The argument does not depend on its content.\par\medskip
\ifdefined\hypertarget\hypertarget{sourcepointer-3}{}\fi
\noindent\textit{Source reference.} Original lemma 067B: sources/stacks-more-morphisms/pinned-original-source.tex, source lines 17272--17272. Cited by the retained statement or proof but not part of the collected spans; supplied only as a navigation pointer. The argument does not depend on its content.\par\medskip
\ifdefined\hypertarget\hypertarget{sourcepointer-4}{}\fi
\noindent\textit{Source reference.} Original lemma 09VA: sources/stacks-more-morphisms/pinned-original-source.tex, source lines 17272--17272. Cited by the retained statement or proof but not part of the collected spans; supplied only as a navigation pointer. The argument does not depend on its content.\par\medskip
\begin{lemma}
\label{lemma-composition-relative-pseudo-coherent}
Let $f : X \to Y$ be a morphism of schemes locally of finite type
over a base $S$. Let $m \in \mathbf{Z}$. Let $E$ be an object of
$D(\mathcal{O}_X)$. Assume $\mathcal{O}_Y$ is pseudo-coherent relative
to $S$\footnote{This means $Y \to S$ is pseudo-coherent, see
Definition \ref{definition-pseudo-coherent}.}.
Then the following are equivalent
\begin{enumerate}
\item $E$ is $m$-pseudo-coherent relative to $Y$, and
\item $E$ is $m$-pseudo-coherent relative to $S$.
\end{enumerate}
\end{lemma}

\begin{proof}
The question is local on $X$, hence we may assume $X$, $Y$, and $S$ are affine.
Arguing as in the proof of More on Algebra, Lemma
\ref{more-algebra-lemma-pull-relative-pseudo-coherent}
we can find a commutative diagram
$$
\xymatrix{
X \ar[r]_i \ar[d]_f &
\mathbf{A}^m_Y \ar[r]_j \ar[ld]^p &
\mathbf{A}^{n + m}_S \ar[ld] \\
Y \ar[r] &
\mathbf{A}^n_S
}
$$
The assumption that $\mathcal{O}_Y$ is pseudo-coherent relative to $S$
implies that $\mathcal{O}_{\mathbf{A}^m_Y}$ is pseudo-coherent relative
to $\mathbf{A}^m_S$ (by flat base change; this can be seen by using
for example Lemma \ref{lemma-base-change-relative-pseudo-coherent}).
This in turn implies that $j_*\mathcal{O}_{\mathbf{A}^n_Y}$ is
pseudo-coherent as an
$\mathcal{O}_{\mathbf{A}^{n + m}_S}$-module. Then the equivalence of
the lemma follows from
Derived Categories of Schemes, Lemma
\ref{perfect-lemma-closed-push-pseudo-coherent}.
\end{proof}

\end{document}
